% Relatório do TCC2: edição mantida diretamente neste arquivo LaTeX.
\chapter{Introdução}\label{cap:tcc2-1}

Neste capítulo são apresentados o contexto do trabalho, o problema abordado e os objetivos do desenvolvimento. Também é delimitado o conjunto de funcionalidades da versão descrita no relatório, considerando as diferenças entre a proposta inicial e o sistema implementado.

\section{Contextualização}\label{sec:1.1}

O desenvolvimento de um jogo envolve a integração entre sua programação e os elementos com os quais o jogador interage. No jogo \textit{Jump'n'Shoot}, desenvolvido neste trabalho, essa integração compreendeu o controle de um personagem, a construção de um cenário tridimensional, a apresentação de informações na interface e o processamento de uma partida compartilhada pela rede. A motivação registrada no TCC1 foi utilizar a criação de um jogo como oportunidade de aprendizado sobre essas diferentes atividades, tendo como proposta um jogo de tiro em primeira pessoa com elementos de plataforma.

Para a implementação foi utilizado o motor de jogos Godot. Sua organização em nós e cenas permite representar tanto elementos do mundo, como personagens e armas, quanto telas e componentes de interface. Cenas podem ser reutilizadas e combinadas para formar sistemas maiores, enquanto sinais permitem comunicar a ocorrência de eventos entre os elementos. Essa organização oferece uma base para reunir, em um mesmo projeto, as diferentes partes necessárias à execução de um jogo. \cite{GODOT-CONCEITOS}

A participação de vários jogadores exige manter informações coerentes entre computadores distintos, apesar dos atrasos de comunicação. No jogo desenvolvido, estados como posição, coleta de itens e pontuação foram tratados por um servidor dedicado, com mensagens de eventos e distribuição periódica de estados da partida. \cite{FIEDLER-REDE}

\section{Problema e motivação}\label{sec:1.2}

O problema deste trabalho consistiu em desenvolver e integrar as funcionalidades de um jogo de tiro em primeira pessoa 3D, competitivo e executado em rede. A implementação deveria permitir que os participantes se movimentassem por uma arena, utilizassem armas, coletassem recursos e disputassem objetivos em equipes. Além das regras de combate, seria necessário organizar a entrada dos jogadores, o início e o encerramento das partidas e a apresentação de seus resultados. Esses requisitos derivaram do documento de projeto elaborado no TCC1.

Uma das principais questões técnicas foi definir qual parte do sistema seria responsável por cada informação. O personagem controlado pelo jogador precisava responder à entrada local, enquanto resultados como dano, consumo de munição e pontuação deveriam ser mantidos no servidor. Essa divisão também exigiu decidir quais estados seriam enviados continuamente, quais poderiam ser comunicados apenas quando mudassem e quais detalhes visuais poderiam permanecer locais. A solução adotada é descrita na seção~\ref{sec:5.2}.

\section{Objetivos}\label{sec:1.3}

\subsection{Objetivo geral}\label{sec:1.3.1}

O objetivo geral do trabalho foi desenvolver um jogo 3D multijogador competitivo do gênero tiro em primeira pessoa, com combate em arena e elementos de plataforma. Na versão implementada, a comunicação foi organizada em uma arquitetura cliente-servidor com servidor dedicado, responsável pelas regras de combate, recursos e resultado das partidas. Esse objetivo retoma a proposta de aprendizado por desenvolvimento do TCC1, incorporando a arquitetura efetivamente construída.

\subsection{Objetivos específicos}\label{sec:1.3.2}

A partir desse objetivo, foram considerados os seguintes objetivos específicos, relacionados à proposta original e às funcionalidades documentadas no projeto.

\begin{itemize}
\item Implementar movimentação em primeira pessoa, corrida, salto duplo e interação com plataformas de impulso.
\item Desenvolver armas, granadas, vida, morte, retorno à partida e coleta de recursos.
\item Implementar modos de jogo competitivos em equipes, com acompanhamento de pontuação e encerramento da partida.
\item Organizar a comunicação entre cliente e servidor e compartilhar contratos e definições usados pelos dois projetos.
\item Construir uma arena e integrar os modelos, materiais, efeitos sonoros e componentes de interface necessários à experiência de jogo.
\item Preparar configurações de exportação para os executáveis de cliente e servidor.
\end{itemize}

\section{Escopo do trabalho}\label{sec:1.4}

O jogo implementado possui dois times, identificados pelas cores azul e vermelho, e dois modos selecionáveis: \textit{Team Deathmatch} (TDM) e \textit{King of the Hill} (KOTH). Ambos utilizam o mapa desenvolvido para o projeto, denominado ``Highrise''. As predefinições de partida estabelecem o limite de 300 segundos, a meta de 20 pontos e dois jogadores por partida. Esses são os parâmetros padrão da versão implementada. A quantidade máxima de salas e o número máximo de jogadores por partida são configuráveis no código. Complementarmente, um modo experimental ampliou a lotação de uma única sala para um ensaio exploratório local, apresentado nas seções~\ref{sec:ensaio-carga-metodo} e~\ref{sec:ensaio-carga-resultados}. O ensaio não estabelece a capacidade do ambiente de hospedagem nem substitui a avaliação humana de jogabilidade.

O conjunto de armas é composto por pistola, submetralhadora e espingarda, além de granadas. Também foram implementados coletáveis de vida e munição e efeitos temporários de dano e cadência de tiro. \textit{Capture the Flag}, canhão laser, lança-foguetes e vidas limitadas fizeram parte da proposta ou do planejamento, mas não compõem o funcionamento da versão implementada. Um mapa extra foi mantido para testes e prototipação dos modelos 3D, sem integrar as predefinições apresentadas no menu principal.

O trabalho contempla a exportação de dois executáveis, um para o cliente, utilizado pelo jogador, e outro para o servidor, para cada um dos sistemas operacionais Windows e Linux. A verificação funcional percorreu o fluxo de partida, conforme descrito na seção~\ref{sec:4.7}. A preparação dos executáveis é descrita na seção~\ref{sec:5.2.11}.

\section{Estrutura do relatório}\label{sec:1.5}

O Capítulo~\ref{cap:tcc2-2} apresenta a fundamentação necessária para caracterizar o jogo. O Capítulo~\ref{cap:tcc2-3} resume as inspirações visuais e de jogabilidade. O Capítulo~\ref{cap:tcc2-4} descreve as ferramentas, a arquitetura e as tecnologias empregadas, o processo de desenvolvimento e os procedimentos de verificação funcional e de ensaio de carga. O Capítulo~\ref{cap:tcc2-5} apresenta primeiro o produto pela perspectiva do jogador e, em seguida, seu desenvolvimento técnico e os resultados disponíveis. Por fim, o Capítulo~\ref{cap:tcc2-6} discute os objetivos alcançados, as contribuições e as possibilidades de continuidade.
